\documentclass{jsarticle}
\usepackage{amsmath}
\usepackage{amssymb}
\begin{document}
\title{}
\author{}
\date{}
\maketitle

   去年に送ったどら焼き、先年に、睦世お姉さんが、銀座からアイスどらやきを
   もらって、夏でした、美味しかったので、福山市はクワイだけでないと、
   虎屋があると言いたくて、後、姉たちが、ドラえもんとは違うが、
   これは、匠生くんが、英語関係の神学校に行ってくれて、柔道までしてくれて、
   俊咲くんが、私がずっと昔に、
   県立図書館で借りた本の内容が、東邦大学の試験にそのまま出題されて、
   そのときに、聞かれていない内容まで、どういう原理の時間の性質かを知らずに、
   そのままの内容を、書いて、石岡先生が、これを知っていたらしく、
   それはだめですよと言われていて、試験日もトキ子おばあちゃんの誕生日で、でも、
   全部、2015年と2016年の剣道で、家族全員の縁で助かって、
   本当に、姉たちがドラえもんで、でも、全員、スーツを着ると、モデルさんであり、
   姉二人に、どら焼きと羊羹、虎屋を送ります。
   1990年の修学旅行のお土産に、山口県萩市の風鈴を、買って帰ったら、
   家族に、”雅くん、選ぶお土産、いいよ”と言われて、今思うと、私、縁起いいか
   と思う次第です。風鈴の鈴も送ります。
   今後も、よろしくお願いいたします。　
   後、匠生くんが、アメルは、高血糖になるから、エチゾラム・ブロチゾラム
   にしてあげて、と、朝ゴハンのフォローをしてくだり、好きです！


\end{document}
